\section{Images and graphics}

\subsection{Excursion on image types}
%\begin{frame}
%	\begin{block}{Color spectra in informatics}
%		
%	\end{block}
%\end{frame}

\begin{frame}
	\begin{columns}[onlytextwidth]
	\column{0.6\textwidth}
		\begin{block}{Bitmap-Images}
			Every image/graphic consists of a specific $length$ and $width$. It is
			separated into a grid of \emph{dots} (or \emph{pixel}):
			Classically, for every dot the corresponding color is stored:
		\end{block}
			
		\begin{exampleblock}{Examples}
			\begin{itemize}
				\item \href{http://de.wikipedia.org/wiki/JPEG}{Joint
					Photographic Experts Group format (JPEG)}
				\item \href{http://de.wikipedia.org/wiki/Portable_Network_Graphics}{Portable Network
					Graphics (PNG)}
				\item \href{http://de.wikipedia.org/wiki/GIF}{Graphics
					Interchange format (GIF)}
			\end{itemize}
		\end{exampleblock}
	\column{0.38\textwidth}
		{\centering\mimage{graphics-dots-explained.pdf}}

		{\tiny derived from \href{http://commons.wikimedia.org/wiki/File:Aufloesung\_erklaert.svg}{CC-BY-SA
		v3.0 by Friedrich Graf}}
	\end{columns}
\end{frame}

\begin{frame}
	\begin{columns}[onlytextwidth]
	\column{0.38\textwidth}
		\mimage{graphics-svg-explained.pdf}
	\column{0.6\textwidth}
		\begin{block}{Vector graphics}
			Store information about elements (e.g. circles and lines) to be drawn, with information
			about the size and relation to each other.
		\end{block}
		\begin{exampleblock}{Examples}
			\begin{itemize}
			\item \href{http://de.wikipedia.org/wiki/SVG}{Scalable Vector Graphics (SVG)}
			\item \href{http://de.wikipedia.org/wiki/Portable_Document_Format}{Portable	Document Format (PDF)}
			\end{itemize}
		\end{exampleblock}
	\end{columns}
\end{frame}

\begin{frame}
	\begin{block}{Comparing bitmap and vector graphics}
		Vector graphics can be scaled to whatever size without loss of
		precision.
	\end{block}

	\begin{columns}[onlytextwidth]
	\column{0.49\textwidth}
		\begin{exampleblock}{PNG image at 5, 50, 100 pixel size}
			\includegraphics[width=5px]{default-image.png}
			\includegraphics[width=50px]{default-image.png}
			\includegraphics[width=100px]{default-image.png}
		\end{exampleblock}
	\column{0.49\textwidth}
		\begin{exampleblock}{PDF image at same sizes}
			\includegraphics[width=5px]{default-image.pdf}
			\includegraphics[width=50px]{default-image.pdf}
			\includegraphics[width=100px]{default-image.pdf}
		\end{exampleblock}
	\end{columns}
\end{frame}
		
\subsection{Including graphical images}
\begin{frame}
	\begin{block}{Using the \lp{graphicx} package}
		Images can be included with the \lcs{includegraphics} command.
		They can be stored at different locations provided by \lcs{graphicspath}.
	\end{block}
	\codex{Simple image inclusion}{includegraphics}
	{\tiny \href{../doc/graphicx.pdf}{(Graphixs-Documentation)}}
\end{frame}
\begin{frame}
	\codex{Including scaled images}{includegraphics-scaled}
\end{frame}
\begin{frame}
	\codex{Including images with defined with and height}{includegraphics-textwidth}
\end{frame}
\begin{frame}
	\codex{Display part of image}{includegraphics-viewport}
\end{frame}
